\section{Microscopic acquisition with the complete population}\label{sec:micro}

\subsection{Freezing one analytic chart}

Use \eqref{eq:band} and the scaled coordinate
\[
 z=a+w/L,\qquad N=\lfloor\ee^{L/2}\rfloor.
\]
Time, cutoff, and all center data are held fixed in every \(w\) or
\(z\) derivative. Define
\begin{align}
 d_n&=n^{-1/2}\exp\left\{\frac{c}{4L}
                       [(\log n)^2-L\log n]\right\},\\
 \lambda_n&=\frac14-\frac{\log n}{2L}\in[0,1/4],\qquad
 \phi_n=\Theta(T,a)-\nu(a)\log n,\label{eq:frozendata}\\
 P_a(w)&=\sum_{n\le N}d_n\cos(\phi_n+\lambda_nw),\\
 K_a(w)&=\frac{\ee^{-\pi^2T/64}\mathcal R_N(a+w/L)}
                         {A(T,a+w/L)},\qquad
 q_a=P_a+K_a,\qquad f_a(w)=\frac{H_T(a+w/L)}{A(T,a+w/L)}.
 \label{eq:fq}
\end{align}
The nonvanishing carrier \(A\) is a normalization, not a lower bound for
the field. Both reflected arms are in the cosine, and \(K_a\) is retained
as a whole analytic function.

The positive coefficient moments supply both the freezing cost and the
later error payments.
\begin{lemma}[Fixed-order coefficient moments]\label{lem:moments}
Let \(\delta=1/4-c/16\in[3/16,7/32]\). For every fixed integer \(k\ge0\),
\begin{equation}\label{eq:moments}
 M_k:=\sum_{n\le N}d_n\lambda_n^k
 =\frac{2k!}{L^k}\ee^{\delta L}\bigl(1+O_k(L^{-1})\bigr).
\end{equation}
For \(k=0\), the convention is \(\lambda_n^0=1\), including a possible
\(\lambda_N=0\).
\end{lemma}
\begin{proof}
Extend the summand to \(1\le x\le\ee^{L/2}\). It is nonnegative and
decreasing, because
\[
 \frac{\dd\log d(x)}{\dd\log x}
 =-\frac12-\frac c4+\frac{c\log x}{2L}\in[-3/4,-1/2]
\]
and \(\lambda(x)\) decreases. Integral comparison, including the final
fractional interval, has discrepancy \(O_k(1)\). The change of variables
\(v=L/2-\log x\) gives the exact integral
\[
 \frac{\ee^{\delta L}}{(2L)^k}
 \int_0^{L/2}v^k\exp\left(-\frac v2+\frac{cv^2}{4L}\right)\dd v.
\]
On this interval the exponential is at most \(\ee^{-3v/8}\), and its
difference from \(\ee^{-v/2}\) is at most
\(v^2\ee^{-3v/8}/(4L)\). The integral is therefore
\[
 \int_0^\infty v^k\ee^{-v/2}\dd v+O_k(L^{-1})
 =2^{k+1}k!+O_k(L^{-1}).
\]
The exponential scale absorbs the \(O_k(1)\) discretization error,
uniformly since \(\delta\ge3/16\).
\end{proof}

\begin{theorem}[Complete local approximation]\label{thm:local}
For every fixed \(R>1\), and for sufficiently large \(a\), the inherited
field input gives
\begin{equation}\label{eq:localerror}
 \sup_{|w|\le R}|\partial_w^j(f_a-q_a)|
       \le C_Ra^{-3/4},\qquad 0\le j\le4.
\end{equation}
Moreover, with \(\chi=a^{-1/4}\ee^{-cL/16}\),
\begin{equation}\label{eq:Ksmall}
 \sup_{|w|\le R}|\partial_w^jK_a(w)|\le C_R\chi,
                    \qquad 0\le j\le4.
\end{equation}
The constants and onsets are finite but not numerically instantiated.
\end{theorem}
\begin{proof}
Work first on \(|w|\le2R\), increasing that fixed radius if necessary
for a Cauchy margin. Eventually this lies inside the declared fixed
provider target. Since
\[
 \Ima\Theta(T,a+w/L)=\Ima w/4+O_R((aL^2)^{-1}),
\]
the ratio \((|C_+|+|C_-|)/|A|\) is bounded by a constant depending on
\(R\). Thus \eqref{eq:provider}, after division by \(A\), is
\(O_R(a^{-3/4})\).

Uniformly for \(n\le N\),
\begin{align*}
 L(a+w/L)-L&=O_R((aL)^{-1}),\\
 w_n(T,a+w/L)/d_n&=1+O_R((aL)^{-1}),\\
 \Theta(T,a+w/L)-\nu(a+w/L)\log n
    &=\phi_n+\lambda_nw+O_R((aL^2)^{-1}).
\end{align*}
The last estimate follows from
\(\Theta_z=L(z)/4+\pi T/(16z)\) and \(\Theta_{zz}=O(1/a)\).
Because \(|\lambda_n w|\le R/2\) on the larger disk, summing the
holomorphic freezing error costs at most
\[
 O_R\!\left(\frac{M_0}{aL}\right)
 =O_R(a^{-25/32}/L)=O_R(a^{-3/4}).
\]
The heat correction in the phase and the quadratic weights have been
frozen at their actual values, not discarded.

Apply Cauchy's estimate to this holomorphic difference on the larger
disk. In particular, if its supremum there is \(C_{2R}a^{-3/4}\),
then at any point with a radius-\(R\) Cauchy circle available the
order-\(j\) bound is \(j!R^{-j}C_{2R}a^{-3/4}\).
This proves \eqref{eq:localerror} for the stated finite set of orders.

Finally the correction input gives
\(\mathcal R_N(z)/A(0,z)=O_R(a^{-1/4})\), and exactly
\[
 \ee^{-\pi^2T/64}\frac{A(0,z)}{A(T,z)}
                 =\ee^{-T L(z)^2/16}.
\]
Its modulus on the disk is \(O_R(\ee^{-cL/16})\).
Another Cauchy estimate proves \eqref{eq:Ksmall}. No moving floor or
real-line big-\(O\) term has been differentiated.
\end{proof}

The inherited terminal population scale and earlier phase-retaining
first-derivative approximation belong to \cite{collision,s04}.
The RSCT applications are the fixed-\(k\) constants in
\eqref{eq:moments}, the event-subtracted acquisition, the particular
signed error below, and the fourth-derivative use in
Section~\ref{sec:mirror}.

\subsection{Subtract the necessary event equations first}

For an analytic function \(g\), write
\[
 g^\circ(w)=g(w)-g(0)-wg'(0).
\]
At any actual event of multiplicity \(m\ge2\), the nonzero carrier gives
\(f_a(0)=f_a'(0)=0\). Set \(S_a=q_a^\circ\). If \(e=f_a-q_a\),
then \(f_a-S_a=e-e(0)-we'(0)\), so Taylor's integral remainder yields
\begin{equation}\label{eq:eventerror}
 |f_a(w)-S_a(w)|\le C_Ra^{-3/4}|w|^2,\qquad
 |f_a'(w)-S_a'(w)|\le C_Ra^{-3/4}|w|.
\end{equation}
It also gives \(q_a(0),q_a'(0)=O(a^{-3/4})\).
These are just two necessary equations, valid for every \(m\ge2\);
the remaining \(m-2\) equations have not been imposed.
Explicitly,
\begin{equation}\label{eq:subtracted}
\begin{split}
 S_a(w)={}&\sum_{n\le N}d_n
 \{\cos(\phi_n+\lambda_nw)-\cos\phi_n
                       +\lambda_nw\sin\phi_n\}\\
 &+K_a(w)-K_a(0)-wK_a'(0).
\end{split}
\end{equation}

\begin{proposition}[Size of the separate absolute allowances]\label{prop:allowance}
For fixed \(\beta>0\), define
\[
 U_\beta=\sum_{n\le N}d_n(\cosh(\beta\lambda_n)-1),\qquad
 V_\beta=\sum_{n\le N}d_n\lambda_n\sinh(\beta\lambda_n).
\]
Then
\begin{equation}\label{eq:UV}
 U_\beta=\frac{2\beta^2}{L^2}\ee^{\delta L}(1+O_\beta(L^{-1})),
 \qquad
 V_\beta=\frac{4\beta}{L^2}\ee^{\delta L}(1+O_\beta(L^{-1})).
\end{equation}
These bound the value and derivative in \eqref{eq:subtracted},
apart from the \(O_\beta(\chi)\) correction.
\end{proposition}
\begin{proof}
At \(w=\ii\beta\), a subtracted cosine equals
\[
 \cos\phi\,(\cosh x-1)+\ii\sin\phi\,(x-\sinh x),
 \qquad x=\beta\lambda_n.
\]
Because \(\cosh x-1\ge\sinh x-x\ge0\), its maximum modulus over a free
real phase is \(\cosh x-1\). The derivative has maximal modulus
\(\lambda_n\sinh x\). Taylor's formula, with
\(0\le x\le\beta/4\), gives
\[
 U_\beta=\frac{\beta^2}{2}M_2+O_\beta(M_4),\qquad
 V_\beta=\beta M_2+O_\beta(M_4).
\]
Insert Lemma~\ref{lem:moments}.
\end{proof}

The allowances diverge uniformly in \(c\). Event subtraction saves
\(L^2\), but a separate absolute estimate still pays the terminal
population. This is a lower bound on the size of a specified upper
allowance. It is not a lower bound on \(|S_a|\), on its signed quadratic
form, or on the actual error. In particular, the independent maximizing
phases need not satisfy the event equations or the canonical phase law.

\subsection{A small acquisition error for the signed expression}

Fix \(y=\beta/L\), and define
\begin{align}
 k_a&=1+\frac{(\beta/L)^2}{4a^2+(\beta/L)^2},\\
 \mathfrak d_a&=2k_a+\frac\beta L
              \Ima\frac{A_z(T,a+\ii\beta/L)}{A(T,a+\ii\beta/L)},\\
 \cD_a(g)&=\mathfrak d_a|g(\ii\beta)|^2
                   +\beta\Ima(g'(\ii\beta)\overline{g(\ii\beta)}).
 \label{eq:D}
\end{align}
The carrier logarithmic derivative is
\begin{equation}\label{eq:bcarrier}
 b(z):=\frac{A_z}{A}=\frac7{4z}-\frac\pi8+\frac{TL(z)}{8z},
 \qquad \mathfrak d_a=2+O_\beta(a^{-2}L^{-2}).
\end{equation}
Since \(H_T'=A_zf_a+ALf_a'\), the normalization is exact:
\begin{equation}\label{eq:physicalD}
 \cD_a(f_a)=\frac{\cB(a,\beta/L)}{|A(T,a+\ii\beta/L)|^2}.
\end{equation}
At an all-real event,
\begin{equation}\label{eq:physicalDnegative}
 \cD_a(f_a)=
 \left[(2-m)k_a-\frac{\beta^2}{L^2}
  \sum_{r\ne\pm a}\frac1{(a-r)^2+\beta^2/L^2}\right]|f_a(\ii\beta)|^2<0.
\end{equation}

\begin{proposition}[Complete signed-error payment]\label{prop:Derror}
Under an actual event, \(\cD_a(f_a)\) differs from
\(\cD_a(S_a)\) by at most
\begin{equation}\label{eq:eta}
 \eta_a=O_\beta\left(
       a^{-3/4}\ee^{\delta L}/L^2+a^{-3/2}\right)
       =o(L^{-2}).
\end{equation}
\end{proposition}
\begin{proof}
If the value and derivative errors at \(\ii\beta\) are
\(\epsilon_0,\epsilon_1\), expansion gives
\begin{align*}
 |\cD_a(f_a)-\cD_a(S_a)|
 \le{}&|\mathfrak d_a|(2|S_a|\epsilon_0+\epsilon_0^2)\\
 &+\beta(|S_a'|\epsilon_0+|S_a|\epsilon_1+\epsilon_0\epsilon_1).
\end{align*}
Use \eqref{eq:eventerror}, Proposition~\ref{prop:allowance}, and
\eqref{eq:Ksmall}. Since \(\delta\le7/32\), the slowest power in the
first term is \(a^{-17/32}/L^2\), proving the claim.
\end{proof}

This is absolute normalized accuracy. It gives neither a nonzero
curvature nor a favorable sign for \(\cD_a(S_a)\). To exclude an event,
one would need a lower bound exceeding \(\eta_a\) for the actual
correlated expression.

Two premises of the outer-line shortcut fail here. The formal
shifted-zeta real part approaches \(1/2+c/4<1\), outside the absolutely
convergent series and Euler product used in Section~\ref{sec:fixed};
and
\[
 |C_+/C_-|(T,a+\ii\beta/L)\longrightarrow\ee^{-\beta/2}.
\]
The reflected arm remains present. These facts delimit that shortcut;
they do not exclude other arithmetic comparisons.
